% ============================================================================
% ARBEITSBLATT: Wiederholung Bewegung (LP-05)
% Physik Klasse 6 - Übungsstunde vor Test
% Stand: 29.01.2026
% ============================================================================

\documentclass[11pt,a4paper]{article}

\usepackage[utf8]{inputenc}
\usepackage[T1]{fontenc}
\usepackage[ngerman]{babel}
\usepackage{geometry}
\usepackage{helvet}
\usepackage{tikz}
\usepackage{enumitem}
\usepackage{fancyhdr}
\usepackage{xcolor}
\usepackage{tcolorbox}
\usepackage{amssymb}
\usepackage{amsmath}
\usepackage{siunitx}
\usepackage{qrcode}
\usepackage[hidelinks]{hyperref}

\sisetup{locale=DE, per-mode=symbol}

\renewcommand{\familydefault}{\sfdefault}

\geometry{left=1.5cm, right=1.5cm, top=1.5cm, bottom=1.2cm}

% Fachfarbe Physik
\definecolor{physik}{HTML}{2c5aa0}
\colorlet{fachfarbe}{physik}

% Hilfsfarben
\definecolor{gridcolor}{HTML}{c0d0e8}

% ============================================================================
% KOPFZEILE
% ============================================================================
\pagestyle{fancy}
\fancyhf{}
\fancyhead[L]{\small\textbf{Physik Klasse 6} | Wiederholung Bewegung}
\fancyhead[R]{\small Arbeitsblatt}
\renewcommand{\headrulewidth}{0.4pt}

% ============================================================================
% MERKKASTEN (weißer Hintergrund, farbiger Rahmen)
% ============================================================================
\tcbuselibrary{skins}

\newtcolorbox{merkkasten}[1][]{
    colback=white,
    colframe=fachfarbe,
    coltitle=black,
    colbacktitle=white,
    fonttitle=\bfseries\small,
    boxrule=1pt,
    arc=2pt,
    left=8pt, right=8pt, top=6pt, bottom=6pt,
    #1
}

% ============================================================================
% AUFGABENTITEL
% ============================================================================
\newcommand{\aufgabe}[1]{\noindent\textbf{\textcolor{fachfarbe}{#1}}}

\begin{document}

% ============================================================================
% TITEL MIT QR-CODE
% ============================================================================
\noindent
\begin{minipage}[t]{0.82\textwidth}
    \vspace{0pt}
    {\Large\bfseries Wiederholung: Bewegung}\\[0.3em]
    {\small Name: \underline{\hspace{5cm}} \hfill Datum: \underline{\hspace{3cm}}}
\end{minipage}
\hfill
\begin{minipage}[t]{0.15\textwidth}
    \vspace{0pt}
    \raggedleft
    \href{https://devbox42.github.io/sciencesim/physik/kl06/02-bewegung/std-05-wiederholung/LP-05-wiederholung.html}{\qrcode[height=1.5cm]{https://devbox42.github.io/sciencesim/physik/kl06/02-bewegung/std-05-wiederholung/LP-05-wiederholung.html}}
\end{minipage}
\vspace{0.2em}
\hrule
\vspace{0.6em}

% ============================================================================
% BLOCK A: BEWEGUNGSARTEN
% ============================================================================
\aufgabe{Aufgabe A: Kreuze die richtige Bahnform und Bewegungsart an.}

\vspace{0.3em}
{\small
\begin{tabular}{|p{5.5cm}|c|c|c|c|}
\hline
\textbf{Beispiel} & \multicolumn{2}{c|}{\textbf{Bahnform}} & \multicolumn{2}{c|}{\textbf{Bewegungsart}} \\
\cline{2-5}
 & geradl. & krumml. & gleichf. & ungleichf. \\
\hline
a) ICE fährt mit 200\,km/h geradeaus & $\square$ & $\square$ & $\square$ & $\square$ \\
\hline
b) Auto bremst an der Ampel & $\square$ & $\square$ & $\square$ & $\square$ \\
\hline
c) Karussell dreht gleichmäßig & $\square$ & $\square$ & $\square$ & $\square$ \\
\hline
d) Fahrrad fährt eine Kurve & $\square$ & $\square$ & $\square$ & $\square$ \\
\hline
e) Rolltreppe fährt gleichmäßig hoch & $\square$ & $\square$ & $\square$ & $\square$ \\
\hline
f) Achterbahn fährt durch Looping & $\square$ & $\square$ & $\square$ & $\square$ \\
\hline
\end{tabular}
}

\vspace{0.6em}

% ============================================================================
% BLOCK B: GESCHWINDIGKEIT
% ============================================================================
\begin{merkkasten}[title=Formel: Geschwindigkeit]
\begin{center}
{\large $v = \dfrac{s}{t}$} \qquad umgestellt: \quad $t = \dfrac{s}{v}$ \qquad $s = v \cdot t$
\end{center}
\end{merkkasten}

\vspace{0.4em}

\aufgabe{Aufgabe B: Berechne die fehlende Größe. Schreibe Formel und Rechnung auf.}

\vspace{0.3em}
{\small
\renewcommand{\arraystretch}{1}
\begin{tabular}{|c|p{8cm}|p{6.5cm}|}
\hline
& \textbf{Angaben} & \textbf{Rechnung} \\
\hline
B1 & Radfahrer: $s = 200$\,m, $t = 40$\,s. Berechne $v$. & \\[1.2cm]
\hline
B2 & Jogger: $v = 4$\,\si{\metre\per\second}, $s = 200$\,m. Berechne $t$. & \\[1.2cm]
\hline
B3 & Skater: $v = 6$\,\si{\metre\per\second}, $t = 30$\,s. Berechne $s$. & \\[1.2cm]
\hline
B4 & Schwimmer: $s = 150$\,m, $t = 30$\,s. Berechne $v$. & \\[1.2cm]
\hline
B5 & Läufer: $v = 8$\,\si{\metre\per\second}, $s = 160$\,m. Berechne $t$. & \\[1.2cm]
\hline
B6 & Bus: $v = 5$\,\si{\metre\per\second}, $t = 40$\,s. Berechne $s$. & \\[1.2cm]
\hline
\end{tabular}
}

\vspace{0.6em}

% ============================================================================
% BLOCK C: s-t-DIAGRAMME
% ============================================================================
\aufgabe{Aufgabe C1: Lies die Werte aus dem Diagramm ab.}

\vspace{0.3em}
\begin{minipage}{0.55\textwidth}
\begin{tikzpicture}[scale=0.65]
    \draw[gridcolor, very thin] (0,0) grid (8,5);
    \draw[->, thick] (0,0) -- (8.5,0) node[right] {$t$ in s};
    \draw[->, thick] (0,0) -- (0,5.5) node[above] {$s$ in m};
    \foreach \x in {0,2,4,6,8} { \node[below] at (\x,0) {\small \x}; }
    \foreach \y/\label in {0/0,1/20,2/40,3/60,4/80} { \node[left] at (0,\y) {\small \label}; }
    \draw[fachfarbe, very thick] (0,0) -- (8,4);
    \draw[fachfarbe, thick] (3.85,1.85) -- (4.15,2.15);
    \draw[fachfarbe, thick] (3.85,2.15) -- (4.15,1.85);
    \draw[fachfarbe, thick] (5.85,2.85) -- (6.15,3.15);
    \draw[fachfarbe, thick] (5.85,3.15) -- (6.15,2.85);
\end{tikzpicture}
\end{minipage}
\hfill
\begin{minipage}{0.42\textwidth}
a) Weg nach $t = 4$\,s: \quad $s$ = \rule{2cm}{0.4pt}

\vspace{0.6em}
b) Zeit bei $s = 60$\,m: \quad $t$ = \rule{2cm}{0.4pt}

\vspace{0.6em}
c) Weg nach $t = 2$\,s: \quad $s$ = \rule{2cm}{0.4pt}

\vspace{0.6em}
d) Zeit bei $s = 80$\,m: \quad $t$ = \rule{2cm}{0.4pt}
\end{minipage}

\newpage

% ============================================================================
% SEITE 2
% ============================================================================

\aufgabe{Aufgabe C2: Vergleiche die Läufer und begründe.}

\vspace{0.5em}
\begin{minipage}{0.5\textwidth}
\begin{tikzpicture}[scale=0.8]
    \draw[gridcolor, very thin] (0,0) grid (8,5);
    \draw[->, thick] (0,0) -- (8.5,0) node[right] {$t$ in s};
    \draw[->, thick] (0,0) -- (0,5.5) node[above] {$s$ in m};
    \foreach \x in {0,2,4,6,8} { \node[below] at (\x,0) {\small \x}; }
    \foreach \y/\label in {0/0,1/20,2/40,3/60,4/80} { \node[left] at (0,\y) {\small \label}; }
    \draw[red, very thick] (0,0) -- (4,4);
    \node[red] at (4.5,4) {\small A};
    \draw[blue, very thick] (0,0) -- (8,4);
    \node[blue] at (8.3,4) {\small B};
\end{tikzpicture}
\end{minipage}
\hfill
\begin{minipage}{0.47\textwidth}
a) Kreuze an: Wer ist schneller?

\vspace{0.3em}
\quad $\square$ Läufer A \quad $\square$ Läufer B

\vspace{1em}
b) Begründe deine Antwort:

\vspace{2.5cm}

c) Berechne $v$ für Läufer A:

\vspace{2cm}

d) Berechne $v$ für Läufer B:

\vspace{2cm}
\end{minipage}

\vspace{1em}

% ============================================================================
% BLOCK D: VERKEHRSSICHERHEIT
% ============================================================================
\begin{merkkasten}[title=Verkehrssicherheit]
\textbf{Anhalteweg} = Reaktionsweg + Bremsweg \quad | \quad Bei 1\,s Reaktionszeit: $s_R$ = $v$
\end{merkkasten}

\vspace{0.4em}

\aufgabe{Aufgabe D: Berechne Reaktionsweg und Anhalteweg. Trage die Werte ein.}

\vspace{0.2em}
{\small
\renewcommand{\arraystretch}{1}
\begin{tabular}{|c|c|c|c|c|}
\hline
& \textbf{Geschwindigkeit} & \textbf{Reaktionsweg} & \textbf{Bremsweg} & \textbf{Anhalteweg} \\
\hline
D1 & 5\,\si{\metre\per\second} & & 2,5\,m & \\[1.2cm]
\hline
D2 & 10\,\si{\metre\per\second} & & 10\,m & \\[1.2cm]
\hline
D3 & 8\,\si{\metre\per\second} & & 6,4\,m & \\[1.2cm]
\hline
D4 & 15\,\si{\metre\per\second} & & 22,5\,m & \\[1.2cm]
\hline
D5 & 12\,\si{\metre\per\second} & & 14,4\,m & \\[1.2cm]
\hline
\end{tabular}
}

\vspace{0.6em}

\aufgabe{Aufgabe E: Erkläre, warum der Sicherheitsabstand bei höherer Geschwindigkeit größer sein muss.}

\vspace{2cm}

\end{document}
